\section{Preliminary Literature Review}
Kernel-bypass and in-kernel fast-path research has established a strong foundation for contemporary cloud-native networking. The XDP model introduced by H\o iland-J\o rgensen \textit{et al.} demonstrated that programmable packet processing at early ingress can substantially reduce processing overhead and improve throughput under realistic traffic loads \cite{xdp2018}. This line of work, combined with sustained progress in eBPF verifier capabilities and tooling, has enabled production-grade policy enforcement and observability in systems such as Cilium \cite{ciliumArchitecture,ciliumHubble}.

In parallel, service mesh design has shifted from per-pod proxies toward shared or waypoint-based mediation models, including Istio Ambient Mesh and ztunnel-based transport indirection \cite{istioAmbient,ambientDesign}. These architectures reduce per-workload proxy replication and alter the reconciliation loop between control-plane intent and data-plane realization. Existing performance discussions frequently emphasize CPU and memory savings, startup improvements, and operational ergonomics, but comparatively fewer studies provide rigorous multi-tenant fairness analyses under controlled noisy-neighbor conditions. The literature therefore suggests a maturing implementation ecosystem but an incomplete evidence base regarding isolation guarantees in shared kernel-level planes.
